% Passa hidelinks ao hyperref ANTES do abntex2 carregá-lo
\PassOptionsToPackage{hidelinks}{hyperref}
\documentclass[
  12pt,
  a4paper,
  oneside,
  article,
]{abntex2}

% ── Encoding e fonte ─────────────────────────────────────────
\usepackage[utf8]{inputenc}
\usepackage[T1]{fontenc}
\usepackage{mathptmx}   % Times New Roman em todo o documento
\usepackage{soul}       % \so{} para letter spacing (compatível com pdfLaTeX)
\usepackage{csquotes}

% Times New Roman também nos títulos
\renewcommand{\ABNTEXchapterfont}{\rmfamily}
\renewcommand{\ABNTEXchapterfontsize}{\normalsize}
\renewcommand{\ABNTEXsectionfont}{\rmfamily}
\renewcommand{\ABNTEXsubsectionfont}{\rmfamily}
\renewcommand{\ABNTEXsubsubsectionfont}{\rmfamily}

% ── biblatex-abnt ────────────────────────────────────────────
\usepackage[
  backend=biber,
  style=abnt,
  extrayear=true,
  repeatfields=true,
  giveninits=true,
]{biblatex}
\addbibresource{referencias.bib}
\renewcommand{\mkbibnamefamily}[1]{\MakeUppercase{#1}}

% ── Pacotes extras ───────────────────────────────────────────
\usepackage{enumitem}
\usepackage{indentfirst}
\usepackage{setspace}
\usepackage{array}

% ── Sumário: formatação uniforme para seções e referências ───
\usepackage{titletoc}
\titlecontents{section}[0pt]{\normalfont}
  {\thecontentslabel\enspace\MakeUppercase}{}
  {\hfill\contentspage}
\titlecontents{subsection}[1.5em]{\normalfont}
  {\thecontentslabel\enspace}{}
  {\hfill\contentspage}
% abntex2 usa chapter* para Referências — igualamos ao estilo das seções
\titlecontents{chapter}[0pt]{\normalfont}
  {\MakeUppercase}{\MakeUppercase}
  {\hfill\contentspage}


% ── Paginação: NBR 14724 → canto superior direito,
%    2 cm da borda superior, 2 cm da borda direita ────────────
\setlength{\headheight}{14.5pt}   % evita aviso do fancyhdr
\setlength{\headsep}{0.72cm}      % distância do cabeçalho ao texto
% Com top=3cm e headsep=0.72cm, o número fica a ~2cm da borda superior
% (3cm de margem − 0.72cm de headsep − altura do cabeçalho ≈ 2cm)

\makepagestyle{abntpaginacao}
\makeevenhead{abntpaginacao}{}{}{\thepage}
\makeoddhead{abntpaginacao}{}{}{\thepage}
\makeheadrule{abntpaginacao}{0pt}{0pt}
\pagestyle{abntpaginacao}

% ── Dados do documento ───────────────────────────────────────
\titulo{Implementação de Carteira de Criptomoedas com Suporte a Pagamentos \textit{offline}}
\autor{Dante Martini Chiarelli Ramacciotti}
\local{Bauru}
\data{Março/2026}
\orientador{Kelton Augusto Pontara da Costa}
\instituicao{%
  Universidade Estadual Paulista ``Júlio de Mesquita Filho''
  \par Faculdade de Ciências
  \par Departamento de Computação
}
\tipotrabalho{Trabalho de Conclusão de Curso}
\preambulo{Trabalho de Conclusão de Curso do Curso de Bacharelado em Ciência da Computação da Universidade Estadual Paulista ``Júlio de Mesquita Filho'', Faculdade de Ciências, campus Bauru}

% ══════════════════════════════════════════════════════════════
\begin{document}
\frenchspacing

% ── CAPA ─────────────────────────────────────────────────────
\begin{capa}
	\center
	\normalfont\normalsize\bfseries
	\so{UNIVERSIDADE ESTADUAL PAULISTA} \\
	\so{``JÚLIO DE MESQUITA FILHO''} \\
	\so{FACULDADE DE CIÊNCIAS} \\
	\so{DEPARTAMENTO DE COMPUTAÇÃO} \\
	\so{BACHARELADO EM CIÊNCIA DA COMPUTAÇÃO}

	\vspace*{\fill}

	\normalfont\normalsize\bfseries
	\so{\MakeUppercase{\imprimirtitulo}}

	\vspace*{\fill}

	% Campos em caixas com bordas (igual ao modelo)
	\begin{flushleft}
		\normalfont\normalsize
		\renewcommand{\arraystretch}{1.8}
		\begin{tabular}{|p{9cm}|p{5.5cm}|}
			\hline
			Nome: \imprimirautor & RA: 221023259 \\
			\hline
		\end{tabular}

		\vspace{7mm}
		\begin{tabular}{|p{7cm}|p{7.5cm}|}
			\hline
			Nome do Orientador: Prof.\ Dr.\ \imprimirorientador \\[2em]
			\hline
		\end{tabular}
	\end{flushleft}

	\vspace*{\fill}

	\normalfont\normalsize
	BAURU - SP \\
	\imprimirdata

\end{capa}

% ── FOLHA DE ROSTO ───────────────────────────────────────────
\begin{folhaderosto}
	\begin{center}
		\normalfont\normalsize
		\MakeUppercase{\imprimirautor}
	\end{center}

	\vspace*{4cm}

	\begin{center}
		\normalfont\normalsize\bfseries
		\MakeUppercase{\imprimirtitulo}
	\end{center}

	\vspace*{2.5cm}

	\hfill
	\begin{minipage}{0.5\textwidth}
		\linespread{1.5}\selectfont
		\normalfont\normalsize
		\imprimirpreambulo
		\par\vspace{0.8em}
		Orientador: Prof. Dr. \imprimirorientador.
	\end{minipage}

	\vspace*{\fill}

	\begin{center}
		\normalfont\normalsize
		\MakeUppercase{\imprimirlocal} \\
		\imprimirdata
	\end{center}
\end{folhaderosto}

% ── FICHA CATALOGRÁFICA ──────────────────────────────────────
\begin{fichacatalografica}
\end{fichacatalografica}

% ── FOLHA DE APROVAÇÃO ───────────────────────────────────────
\begin{folhadeaprovacao}
\end{folhadeaprovacao}

% ── AGRADECIMENTOS ───────────────────────────────────────────
\begin{agradecimentos}
\end{agradecimentos}

% ── EPÍGRAFE ─────────────────────────────────────────────────
\begin{epigrafe}
\end{epigrafe}

% ── RESUMO ───────────────────────────────────────────────────
\begin{resumo}
\end{resumo}

% ── ABSTRACT ─────────────────────────────────────────────────
\begin{resumo}[english]
\end{resumo}

% ── LISTAS ───────────────────────────────────────────────────
\listoffigures*
\listoftables*

% ── LISTA DE ABREVIATURAS E SIGLAS ───────────────────────────
\begin{siglas}
\end{siglas}

% ── SUMÁRIO ──────────────────────────────────────────────────
\pdfbookmark[0]{\contentsname}{toc}
\renewcommand{\contentsname}{SUMÁRIO}
\tableofcontents*
\cleardoublepage

% ── TEXTO ────────────────────────────────────────────────────
\textual
\pagestyle{abntpaginacao}

% ══════════════════════════════════════════════════════════════
\section{Introdução}

% ══════════════════════════════════════════════════════════════
\section{Fundamentação Teórica}

\subsection{Revisão bibliográfica}

Alguns trabalhos já foram realizados explorando a possibilidade de pagamentos \textit{offline} no contexto de criptomoedas. O primeiro deles a teorizar pagamentos entre dois usuários estando ambos \textit{offline} foi proposto por \textcite{dmitrienko2017}. Outros de nota, mais focados no contexto de $CBDC$s, incluem os trabalhos de \textcite{videira2023} e \textcite{christodorescu2020}. Um protocolo de transação \textit{offline} utilizando Zero Knowledge Proofs foi estabelecido por \textcite{huayou2026}, mas o ecossistema em si não foi definido nesse trabalho. \textcite{ianov2021} propuseram um sistema utilizando blockchains pública e privada, mas que requer que o \textit{sender} esteja \textit{online} no momento do pagamento, servindo de \textit{proxy} para conexões de um terminal de pagamento. \textcite{jie2024} propõem uma implementação com prova formal de segurança, embora algumas premissas adotadas não possam ser sempre observadas.

As abordagens citadas aqui serão discutidas com mais detalhes na seção seguinte.

\subsection{Análise das abordagens existentes}

\subsubsection{Dmitrienko et al.}

O artigo foi a primeira solução proposta de pagamentos \textit{offline} em criptomoedas. Os principais pontos abordados incluem o uso de hardware seguro (paralelo às carteiras \textit{tamper-proof} usadas nesse trabalho), a verificação adiada de parâmetros e o emprego de mecanismos probabilísticos de checagem de \textit{timestamps} para, juntamente a um limite de valor de transação, inviabilizar práticas de \textit{double spending}. Estas, somente possíveis em uma eventual quebra do sigilo do $TEE$, se tornariam economicamente inviáveis.

Um dos problemas dessa abordagem é que a checagem feita por meio dos \textit{timestamps} é apenas difícil de ser burlada, não criptograficamente impossível. O trabalho atual tenta prover um método criptograficamente impossível de ser burlado, salvo pela quebra do sigilo do $TEE$.

\subsubsection{Videira}

O artigo apresenta uma abordagem para moedas $CBDC$, que presumem que uma autoridade central irá aprovar as carteiras e as moedas nas carteiras. O protocolo de aprovação é bem elaborado e inclui a definição de uma blockchain local que poderia ser reaproveitada em alguns casos.

A desvantagem da abordagem é a exigência do apoio de um sistema legal que identifique os participantes de cada operação, pois em um dos casos (\textit{secret.nonce\#} is not valid, seção 2.3.8), se o \textit{receiver} da moeda não colaborar com a falha na operação, ele ficará indevidamente com a moeda. Um sistema legal poderia forçá-lo a reverter a operação, mas baseando-se somente na criptografia (o que é o caso quando se trabalha com uma criptomoeda descentralizada) o protocolo não é infalível.

\subsubsection{Si et al.}

O artigo define um protocolo de verificação de soma zero, onde a validade de uma transação seria comprovada pela equação mudança + valor recebido = valor de entrada, sem revelar os valores originais e o saldo de cada carteira. Alguns fatores de \textit{decoy} são adicionados à equação original para camuflar os valores utilizados e garantir a privacidade dos usuários.

É explicitada a dependência de um \textit{Payment Hub} que autoriza os fundos das carteiras, mas não é elaborado como este funcionaria.

Como a preservação da privacidade dos usuários não foi um fator orientador do desenvolvimento da criptomoeda deste presente trabalho, os benefícios da anonimidade \textit{offline} seriam imediatamente desfeitos ao se reconectar novamente à rede.

\subsubsection{Christodorescu et al.}

O artigo propõe o $OPS$ (\textit{Offline Payment System}), um sistema que permite transações entre dois usuários \textit{offline} para moedas $CBDC$, utilizando também $TEEs$. O $OPS$ presume uma hierarquia de autoridades certificadoras ($CAs$) de carteiras semelhantes às $CAs$ que regem os registros $TLS$ que autorizariam fundos \textit{offline}, e essas autoridades seriam por sua vez autenticadas pela autoridade central emitidora da moeda.

A estrutura geral do sistema se assemelha à desenvolvida no presente trabalho, com a diferença mais óbvia sendo que, em vez de uma autoridade central, os fundos seriam autorizados pela rede da criptomoeda utilizando $DTSS$.

\subsubsection{Jie et al.}

O artigo garante defesa contra ataques de \textit{Double Spend}, \textit{Man-in-the-Middle}, \textit{Coin Forgery}, \textit{Offline Transaction Data Forgery} e \textit{Double Deposito}. Ele faz isso assumindo, entre outras coisas, que a comunicação entre usuários é perfeita, o que não pode ser garantido no caso de um atacante malicioso.

Detectei um problema no algoritmo ao não lidar com um caso onde o \textit{receiver} não compartilhava a confirmação da transação, o que foi deixado como \textit{gap} no escopo do artigo.

\subsection{Elaboração da teoria de pagamentos \textit{offline}}

Salvo se forem encontradas vulnerabilidades ou otimizações específicas na teoria, esta etapa está concluída. A teoria, em sua forma atual, é como definida a seguir:

\subsubsection{Terminologia}

Uma função de assinatura e a assinatura resultante são definidas como:
$$\sigma(sk_{signer}, data) = \sigma_{signer}(data)$$

Mensagens são dadas como:
$$ M_{contents} $$

\textit{Dynamic Threshold Signature Scheme} ($DTSS$), ou Esquema de assinatura com limiar dinâmico, em tradução livre:
Um $DTSS$ é um protocolo de distribuição de chave de assinatura entre $n$ participantes de forma que apenas se $t$ ou mais deles colaborem será possível criar uma assinatura. Denomina-se tal sistema um sistema $(t, n)$-threshold.

\subsubsection{Pressupostos}

Esse trabalho foi construído em cima de alguns pressupostos. Estes são:

\begin{enumerate}
	\item As carteiras operando em contextos \textit{offline} têm uma forma segura de comunicação entre si, sem nenhum intermediário malicioso;
	\item Apesar de estudos revelarem possibilidades de violá-los \cite{skorobogatov2012}, os $TEEs$ são invioláveis ou que violá-los seria ineficiente demais para que alguém o fizesse;
	\item As operações na blockchain \textit{online} são confiáveis;
	\item Usuários operam em ambientes com $TEEs$;
	\item As primitivas criptográficas são seguras;
\end{enumerate}

\subsubsection{Modelagem de ameaças}

No protocolo apresentado, supõe-se que um atacante seja capaz de realizar falsificação de moedas, adulteração de dados de transações \textit{offline}, ataques de \textit{Double Spending} e ataques de repetição (\textit{replay}). O atacante pode ser um participante desonesto que adota esses comportamentos para obter ganhos ilícitos.

\paragraph{Falsificação de Moedas}
Um usuário malicioso pode tentar forjar uma $\sigma_{rede}(tx_d)$ sem ter realizado um depósito legítimo, ou inflar o valor dos $UTXOs$ autorizados em sua carteira \textit{offline}.

\paragraph{Adulteração de dados de transações \textit{offline}}

O \textit{sender} ou o \textit{receiver} podem tentar adulterar os dados armazenados no $ANC$, falsificando o histórico de transações da blockchain local.

\paragraph{Ataques de \textit{Double Spending}}
Um usuário malicioso pode tentar gastar os mesmos $UTXOs$ múltiplas vezes em um contexto \textit{offline}, seja com o mesmo \textit{receiver} ou com \textit{receivers} distintos.

\paragraph{Ataques de repetição}
Um atacante pode tentar reutilizar uma transação já registrada na blockchain local para enganar um \textit{receiver} distinto.

\paragraph{Ataques \textit{Sybil} à rede}
Um atacante pode tentar criar \textit{nodes} suficientes para acumular $t$ partes da $CAR$ e forjar assinaturas de depósito.

\subsubsection{A criptomoeda}

A implementação da criptomoeda, com a exceção de sua parte de pagamentos \textit{offline}, terá a mesma base teórica do Bitcoin \cite{nakamoto2008} \cite{bitcoinorg}, com algumas simplificações para que seja possível completar o projeto em tempo. Serão feitas algumas adições para possibilitar os pagamentos \textit{offline}. A mais importante delas é a ideia de que, a partir do momento em que houver na rede mais de $N$ \textit{nodes} conectados que cumpram requisitos específicos (possivelmente ter minerado um número mínimo de blocos), será distribuída, por meio de um protocolo de $DTSS$, uma chave criptográfica de assinatura pela rede, com cada \textit{node} que já foi responsável por consolidar um bloco tendo acesso a uma parte da chave. Outra adição necessária será a de um novo campo booleano em transações para assinalar se estas são depósitos a uma carteira \textit{offline} ou não. O campo, quando mencionado, será representado por $tx.deposit$, sendo $tx$ a transação sendo realizada.

Se a distribuição da chave fosse realizada desde o primeiro momento, isso tornaria a rede vulnerável a um ataque \textit{Sybil}. Aguardar a rede ter mais de $N$ \textit{nodes} conectados que cumpram os requisitos mitiga essa possibilidade, dado que um atacante teria que investir bastantes recursos criando diversos \textit{nodes} e esperando que eles consolidassem blocos, e eventualmente o número de blocos requeridos poderia aumentar, desestimulando um atacante. Essa chave compartilhada entre os \textit{nodes} será chamada de chave de assinatura da rede ($CAR$), enquanto a sua chave de verificação será chamada de chave de verificação da rede ($CVR$). O tamanho específico de $t$ nesse $DTSS$ será explorado mais tarde.

A $CAR$ será a chave utilizada para autorizar transações ($tx$) realizadas para carteiras \textit{offline}, provendo para elas uma assinatura $\sigma_{tx}$. No que será chamado de operação de depósito, cada carteira \textit{offline} irá conter uma cópia da $CVR$ para verificar se os fundos realmente foram autorizados pela rede. Uma vez que uma carteira tem fundos ($UTXOs$) autorizados a serem usados \textit{offline}, esses $UTXOs$ serão congelados na rede, ou seja, não será possível ao seu dono utilizá-los normalmente até que ele próprio solicite um descongelamento e preste contas sobre as transações realizadas. Nesse momento, as transações serão oficializadas na rede, a menos que seus recipientes já tenham oficializado-as previamente.

\subsubsection{A carteira \textit{offline}}

Para evitar tampering enquanto os envolvidos em uma transação estejam \textit{offline}, serão utilizadas carteiras com \textit{Trusted Execution Environments} ($TEEs$), que conterão um registro das operações realizadas. Essas $TEEs$ buscam prevenir que um agente malicioso interfira no funcionamento esperado, e além disso provêem ferramentas que permitem a checagem da validade do ambiente $TEE$, de forma que não seja possível para um atacante emular um ambiente semelhante a fim de enganar o protocolo. Dentro dessas $TEEs$, haveria uma chave de assinatura \textit{offline} ($osk$) que pode autorizar transações (a $osk$ tem como contrapartida a chave de verificação $ovk$). Essa chave seria criada por uma operação dentro da própria $TEE$, de forma que apenas a invasão física do ambiente pudesse revelá-la. As $TEEs$ também realizam o \textit{hashing} do próprio software. Desse modo, quaisquer adulterações no ambiente seriam detectadas e impediriam que a transação acontecesse em primeiro lugar, ou acusariam a adulteração quando o agente malicioso se conectasse novamente à rede para tentar descongelar suas moedas, salvaguardando-as para quando os possíveis recipientes se conectassem e reivindicassem seus $UTXOs$. Isso implica que mesmo se, por algum propósito não malicioso, algo acontecer ao hardware ou ao software, os fundos serão perdidos, mas no universo das criptomoedas esse risco é familiar aos usuários, uma vez que a perda de um dispositivo com determinados tipos de carteira pode ter a mesma implicação.

A carteira será composta pelo $TEE$, que conterá um \textit{counter} $\theta$ de operações realizadas, pelo par de chaves da carteira $(osk, ovk)$, por uma cópia da $CVR$, para uso nas verificações de moedas alheias, e por um ambiente não controlado e não confiável, denominado $ANC$. Este irá armazenar uma blockchain local, que armazenará todas as operações realizadas \textit{offline} assim como todas as operações de depósito. Essa blockchain local também contará com um campo especial, uma assinatura $\sigma(bloco)$, que irá validar se os blocos estão de acordo com as operações que a $TEE$ realizou.

\subsubsection{A autorização de moedas \textit{offline}}
Para que uma carteira possa realizar transações \textit{offline}, é necessário que ela antes realize uma operação de depósito na rede, que irá congelar os $UTXOs$ para uso \textit{offline} e assiná-los para que mais tarde eles sejam checados na verificação da transação \textit{offline}. O depósito ocorre da seguinte forma:

\begin{enumerate}
	\item O usuário submete à rede uma solicitação de depósito a uma carteira \textit{offline} $tx_d$. Essa solicitação de depósito nada mais é do que uma transação com seu campo $tx.deposit$ marcado como \texttt{true}, e o destinatário sendo a $ovk$ de sua carteira \textit{offline}.
	\item A rede verifica a transação como qualquer outra, e uma vez que seja válida, marca os $UTXOs$ como congelados. Se uma nova transação for feita com eles enquanto a carteira \textit{offline} não houver se conectado à rede para fazer a ressincronização, a transação será imediatamente rejeitada. A ressincronização será explicada mais abaixo.
	\item A rede esperará a consolidação de um número $B$ de blocos para minimizar a chance de problemas por conta de \textit{forks} de cadeias, e depois disso criará uma transação de valor nulo para a $ovk$ da carteira \textit{offline} que recebeu os fundos partindo da $CVR$ da rede, contendo $\sigma_{rede}(tx_d)$.
	\item A carteira \textit{offline} faz o download dessa assinatura e a armazena junto da transação em sua blockchain local.
\end{enumerate}

\subsubsection{O protocolo para transações \textit{offline}}

\begin{enumerate}
	\item O \textit{receiver} envia um \textit{nonce} $\delta$ para o \textit{sender}.
	\item O \textit{sender} envia uma transação assinada equivalente a uma transação realizada em um contexto \textit{online}. Além disso, ele envia a assinatura $\sigma_{tx}$ que ele recebe da rede que autenticou seus fundos, uma cópia de sua blockchain, o contador de transações da $TEE$, assim como uma assinatura $\sigma_{osk_{sender}}(\theta_{sender}, \delta) = \sigma_{sender} (\theta_{sender}, \delta)$ sobre o contador e o \textit{nonce}.
	\item A partir desse momento, os fundos utilizados para a operação serão registrados pela $TEE$ como congelados até que o \textit{sender} receba $M_{sucesso}$ ou $M_{erro}$.
	\item O \textit{receiver} irá checar os blocos da blockchain, conferindo se as informações são válidas e se as moedas sendo utilizadas foram em algum momento assinadas pela $CAR$. Ele também deve checar se o número de blocos é fiel ao número $\theta$ recebido. A verificação pode resultar em um sucesso ou em um fracasso.
	      \begin{enumerate}
		      \item Em caso de sucesso, o \textit{receiver} irá registrar a transação na própria blockchain local, assinando-a com sua $osk$ e incrementando o $\theta_{receiver}$ em 1. Ele também irá enviar $M_{sucesso}$ e $\sigma_{receiver}(M_{sucesso})$ para o \textit{sender}.
		      \item Em caso de fracasso, o \textit{receiver} irá registrar em sua blockchain um novo bloco que registra a transação como falha, aumentando o contador de qualquer forma. Ele também irá enviar $M_{erro}$ e $\sigma_{receiver}(M_{erro})$ para o \textit{sender}.
	      \end{enumerate}
\end{enumerate}

\subsubsection{A ressincronização}

Quando a carteira \textit{offline} se conecta à rede, ela deve prestar contas e submeter sua blockchain local para ressincronização. O processo ocorre da seguinte forma:

\begin{enumerate}
	\item A carteira envia sua blockchain local completa à rede.
	\item A rede confere a integridade da cadeia, conferindo as assinaturas $\sigma_{carteira}(bloco)$ criadas pela $TEE$ em cada bloco, garantindo que não houve quaisquer adulterações advindas do $ANC$.
	\item Para cada transação na blockchain local:
	      \begin{enumerate}
		      \item Se o \textit{receiver} ainda não houver submetido à rede, ela é oficializada agora.
		      \item Se o \textit{receiver} já houver submetido, a rede confere a consistência entre as versões.
		      \item Se houver evidência de adulteração, os fundos permanecerão retidos até que o \textit{receiver} se conecte à rede para esclarecer a operação.
	      \end{enumerate}
	\item Os $UTXOs$ não gastos durante o período podem ser reautorizados (se a rede impuser limites de tempo para o gasto de $UTXOs$ \textit{offline}; isso ainda não foi definido) ou transferidos para outra carteira para uso \textit{online}.
\end{enumerate}

Erros possíveis no protocolo:

\begin{enumerate}
	\item \textit{Receiver} recebe $M_{sucesso}$ do $TEE$, a envia para o \textit{sender}, mas não registra a operação:
	      \subitem Isto implicaria na perda do saldo para o \textit{receiver} até que o \textit{sender} se conectasse à rede e submetesse seu histórico de transações.
	\item \textit{Receiver} recebe $M_{sucesso}$ do $TEE$, \textbf{não} a envia para o \textit{sender}, e não registra a operação:
	      \subitem Isso implica na perda dos fundos para ambos. Os fundos estão congelados para o \textit{sender}, então mesmo quando ele se conectar à rede, não poderá recuperá-los. Por outro lado, o \textit{receiver} não tem mais prova da ocorrência da operação. A única forma de recuperá-los seria repetindo o envio detalhado em 2.
	\item \textit{Sender} recebe $M_{sucesso}$ ou $M_{erro}$ mas ignora (seja por má-fé ou por erro de conexão):
	      \subitem Os fundos utilizados na transação estão congelados pela $TEE$ até que uma das mensagens seja recebida. Não há qualquer benefício nessa linha de ação para o \textit{sender}.

	\item \textit{Receiver} recebe $M_{sucesso}$ e não a envia ao \textit{sender}:
	      \subitem Isso seria inconveniente ao \textit{sender}, pois a transação ficaria em aberto, mas não acarretaria em nada além dessa inconveniência. O dinheiro já pertence ao \textit{receiver}, e os fundos estão congelados para o \textit{sender}. Da próxima vez que o \textit{sender} se conectar à internet após o \textit{receiver} ter atualizado a blockchain com suas informações, o mistério da transação em aberto será elucidado.

	\item \textit{Receiver} recebe $M_{erro}$ e não a envia ao \textit{sender}:
	      \subitem Essa é a única possibilidade que permite uma atitude maliciosa. Como os fundos do \textit{sender} estão congelados e só poderão ser utilizados novamente em caso de $M_{erro}$, esses fundos ficarão para sempre inutilizáveis. Por outro lado, supondo um \textit{receiver} honesto e um canal de troca de mensagens sólido, um $M_{erro}$ significaria que o \textit{sender} estaria agindo de má-fé, então os fundos perdidos seriam uma espécie de punição.
\end{enumerate}

% ══════════════════════════════════════════════════════════════
\section{Metodologia}

\subsection{Implementação da criptomoeda}

A implementação da criptomoeda em si está prosseguindo como planejado. Grande parte da implementação já foi concluída. Já foram concluídas:

\begin{itemize}[leftmargin=1.5cm, itemsep=0pt, topsep=2pt]
	\item Estrutura central, com a lógica das carteiras, transações, blocos e cadeia implementadas. A parte criptográfica, de assinaturas e verificações, foi feita usando a \textit{crate} \textit{openssl} \cite{openssl}.
	      \subitem Sistema de seleção de moedas como definido em \textcite{murchbnb}
	      \subitem Sistema customizado de seleção de moedas (não exato).
	      \subitem Algoritmo para cálculo da Merkle Tree (semelhante ao do Bitcoin \cite{nakamoto2008}, mas diferencia-se em como lida com os espaços vazios no nível mais inferior da árvore binária).
	\item Criação de sistema de gerenciamento de threads (\textit{Thread Pool}).
	\item Protocolo HTTP para comunicação \textit{node}—\textit{node} e usuário—\textit{node}.
	\item Criação de interface de terminal (\textit{TUI}) utilizando a \textit{crate} \textit{ratatui}  \cite{ratatui} para gerenciamento de \textit{nodes}.
	\item Implementação de \textit{logger}.
	\item Endpoints:
	      \subitem "/" => \textit{TUI}.
	      \subitem "/status" => Consulta via API de informações sobre o \textit{node}.
	      \subitem "/submit-transaction" => Usuários submetem por aqui transações serializadas para JSON.
	      \subitem "/messages" => Endpoint da API para comunicação entre \textit{nodes}. Suporta mensagens notificando sobre novos blocos ou transações, notificações sobre novos \textit{nodes} encontrados na rede, pedidos por dados (transações ou blocos) e notificações sobre redistribuição da $CAR$.
	\item Sistema de mensagens entre \textit{nodes}.
\end{itemize}

Ainda devem ser elaborados:
\begin{itemize}[leftmargin=1.5cm, itemsep=0pt, topsep=2pt]
	\item Testes para o sistema de mensagens
	\item Resolução de conflitos de cadeia
	\item Sistema de mineração
	\item Sistema de descoberta de \textit{nodes} (servidor que provê os endereços dos \textit{nodes} conhecidos)
	\item Algoritmo de contagem de \textit{nodes} na rede
\end{itemize}

\subsection{Implementação do sistema \textit{offline}}

A implementação do sistema \textit{offline} está estreitamente relacionada com a implementação do resto da criptomoeda, visto que algumas acomodações já são feitas previamente para comportar essa parte do sistema, portanto considero que essa seção está em andamento.

Para concluir o sistema de pagamento \textit{offline}, são necessários apenas dois passos. O primeiro é o protocolo $DTSS$, que deve ser apoiado pelo algoritmo de contagem de \textit{nodes} para poder determinar o $n$ e $t$, sendo $n$ o número de \textit{nodes} ativos que já consolidaram um bloco na cadeia vigente (em contraste com cadeias que podem ter sido descartadas em algum conflito. Estas não deverão ser utilizadas nesse cálculo) e $t$ o número mínimo de \textit{nodes} que devem coordenar para performar a operação de assinatura (denomina-se tal sistema um sistema $(t, n)$-threshold).

Restam a implementação, então, de:

\begin{itemize}[leftmargin=1.5cm, itemsep=0pt, topsep=2pt]
	\item Protocolo de $DTSS$
	\item Assinaturas de carteiras \textit{offline}
	\item Testes de integridade de transações \textit{offline}
\end{itemize}

\subsection{Exploração da carteira \textit{tamper-proof}}

O panorama geral da carteira \textit{tamper-proof} para transações \textit{offline} já foi apresentada em seções anteriores, mas, se possível, esse projeto tem a intenção de investigar com mais profundidade como exatamente esta poderia ser executada, investigando o padrão $TEE$ e a operação de \textit{sealing}, e como adaptar as ferramentas disponíveis para esse uso específico.

\subsection{Exploração dos limiares do sistema $DTSS$}

TODO

\subsection{Testes e avaliação}

Por fim, com a criptomoeda plenamente implementada, será comparada a abordagem desenvolvida ao longo do projeto com outras formas de pagamento \textit{offline} de criptomoedas. Será dado enfoque parcial às teorias, pois as aplicações de software desenvolvidas nesse e em outros projetos tendem a ser protótipos, elaborados com menos atenção do que necessário quando se trabalha com aplicações criptográficas.

\subsection{Redação da monografia}

A partir do próximo mês, iniciará-se a redação da monografia.

\subsection{Cronograma}

\begin{table}[htb]
	\caption{\label{tab:cronograma}Cronograma de atividades}
	\renewcommand{\arraystretch}{1.4}
	\begin{tabular}{|l|c|c|c|c|c|c|c|c|}
		\hline
		\textbf{Atividades}
		                                               & \textbf{Mar} & \textbf{Abr} & \textbf{Mai} & \textbf{Jun}
		                                               & \textbf{Jul} & \textbf{Ago} & \textbf{Set} & \textbf{Out}                 \\
		\hline
		Revisão bibliográfica                          & X            & X            &              &              &   &   &   &   \\
		\hline
		Análise de abordagens existentes               & X            & X            & X            &              &   &   &   &   \\
		\hline
		Elaboração da teoria de pag.\ \textit{offline} & X            & X            & X            &              &   &   &   &   \\
		\hline
		Implementação da criptomoeda                   &              &              & X            & X            & X &   &   &   \\
		\hline
		Implementação do sistema \textit{offline}      &              &              & X            & X            & X &   &   &   \\
		\hline
		Exploração de cart.\ \textit{tamper-proof}     &              &              &              & X            & X &   &   &   \\
		\hline
		Testes e avaliação                             &              &              &              &              & X & X &   &   \\
		\hline
		Redação da Monografia                          &              &              &              &              & X & X & X & X \\
		\hline
	\end{tabular}
\end{table}

% ══════════════════════════════════════════════════════════════
\section{Resultados}

% ══════════════════════════════════════════════════════════════
\section{Conclusão}

% ── REFERÊNCIAS ──────────────────────────────────────────────
\newpage
\postextual
\printbibliography[title=REFERÊNCIAS]

\end{document}
